% CV publico del portafolio (EN). Copia de empleo-agente/salidas/aplicaciones/_cv_base/en/cv.tex.
% Sin nombre de la empresa, sin telefono y sin el abrepuertas (es publico en la web).
% Sin guiones largos (---) a proposito; "--" solo en rangos de fechas.
\documentclass[a4paper,10pt]{article}
\usepackage[top=0.6in, bottom=0.6in, left=0.7in, right=0.7in]{geometry}
\usepackage[T1]{fontenc}
\usepackage[utf8]{inputenc}
\usepackage{lmodern}
\usepackage{microtype}
\usepackage{enumitem}
\usepackage{needspace}
\usepackage{parskip}
\usepackage{titlesec}
\usepackage[dvipsnames]{xcolor}
\definecolor{acc}{HTML}{1F3A5F}
\usepackage[colorlinks=true,urlcolor=acc,linkcolor=acc,bookmarks=false,
  pdfauthor={Juan Nicolás García Vega},pdftitle={CV Juan Nicolás García Vega}]{hyperref}

\pagestyle{empty}
\raggedright
\hyphenpenalty=10000 \exhyphenpenalty=10000
\emergencystretch=3em
\clubpenalty=10000 \widowpenalty=10000

\setlist[itemize,1]{leftmargin=1em,   itemsep=3pt,   parsep=0pt, topsep=2pt,   label=\textbullet}
\setlist[itemize,2]{leftmargin=1.4em, itemsep=1.5pt, parsep=0pt, topsep=1.5pt, label=\textendash}
\titleformat{\section}{\large\bfseries\uppercase}{}{0em}{}[\vspace{0.1em}\titlerule]
\titlespacing{\section}{0pt}{0.8em}{0.5em}
\newcommand{\entry}[3]{\textbf{#1}\ifx&#2&\else, \textit{#2}\fi \hfill \mbox{\textit{#3}}}
\newcommand{\gh}[1]{\href{https://github.com/Nicolas6879/#1}{github.com/Nicolas6879/#1}}

\begin{document}

\begin{center}
    {\huge \textbf{Juan Nicolás García Vega}} \\[0.3em]
    Full-Stack Engineer \textbar{} AI Agents \& Web3 \\[0.35em]
    Bucaramanga, Colombia \textperiodcentered{} Open to relocation or remote \\
    \href{mailto:jngv08@gmail.com}{jngv08@gmail.com} \textperiodcentered{}
    \href{https://www.linkedin.com/in/juan-nicolas-garcia-vega}{linkedin.com/in/juan-nicolas-garcia-vega} \textperiodcentered{}
    \href{https://github.com/Nicolas6879}{github.com/Nicolas6879}
\end{center}

\needspace{6\baselineskip}
\section*{Summary}
Full-stack engineer (B.S. Systems Engineering, Jul.\ 2026) who has put \textbf{8 systems
into production} at an industrial company in Colombia, including an \textbf{ISO 9001 quality system that replaced
600+ paper documents}, and who maintains the company \textbf{ERP} (about 100 users). Builds
AI agents and tools (MCP, a Claude Code plugin). \textbf{Hedera Developer \& Community
Ambassador for Colombia}; \textbf{1st place} Web3 CTF at Hedera DevDay and \textbf{2nd place}
on-chain automation bounty, ETHDenver 2026.

\needspace{6\baselineskip}
\section*{Professional Experience}
\begin{itemize}
    \item \entry{Software Engineer}{Industrial company, Colombia}{Apr.\ 2026 -- Present}
        \begin{itemize}
            \item Lead the move to \textbf{paperless operations} end to end: I design and build
                the systems, including the digital-records app for scanned, filed and signed
                PDFs. Also restructured and completed the earlier attachments system (PDFs
                generated by the ERP, previously handled by hand).
            \item Built an \textbf{MCP server} that lets staff search the QMS catalogs and
                \textbf{publish quality documents from Claude}, so people stop skipping the
                upload step.
            \item Maintain and extend the \textbf{ERP} (Visual Basic, SQL Server), used by about
                \textbf{100 people}.
            \item Part of the team building an internal \textbf{AI assistant} (Gemini API, moving
                to Claude) for company knowledge and reports.
            \item Shipped \textbf{Preoperacionales} (equipment checklists in 3 branches, offline
                mode with auto-sync), \textbf{ComidApp} (multi-plant orders and costing), a
                \textbf{Power BI} dashboard hub, a third-party due-diligence form for
                \textbf{SAGRILAFT} (Colombia's anti-money-laundering rules) with PDF output, and
                a full migration of a \textbf{C\#} system to Vue 3,
                Laravel and SQL Server.
        \end{itemize}

    \item \entry{Systems Engineering Intern}{Industrial company, Colombia}{Aug.\ 2025 -- Apr.\ 2026}
        \begin{itemize}
            \item Built the \textbf{ISO 9001 Quality Management System} from scratch
                (undergraduate thesis), used daily by \textbf{3 departments}, replacing
                \textbf{600+ paper documents}. Six modules: documents, KPIs, audits, risks,
                CAPA and non-conformities, change management. Vue 3, TypeScript, Pinia,
                Laravel, MySQL.
        \end{itemize}

    \item \entry{Developer \& Community Ambassador, Colombia}{Hedera}{2025 -- Present}
        \begin{itemize}
            \item Organized, spoke at and judged \textit{Descifra Hedera: Talk and Web3
                Hackathon} (29 Sep.\ 2026, U26 Fest, Universidad Industrial de Santander):
                3 hours, about \textbf{90 attendees}.
            \item Ran two virtual workshops: the local kickoff of the \textbf{Hedera APEX
                Hackathon} for Mexico (26 Feb.\ 2026, 3 h), and onboarding of students from two
                schools in underserved communities in Brazil to the Hedera Developer Portal,
                to build their first dapp (27 Feb.\ 2026, 3.5 h, part of an inclusion program).
            \item Speaker at \textbf{Blockchain Summit Latam Bogot\'a} (4 Nov.\ 2026): 1-hour
                workshop \textit{Asset Tokenization and AI Agents with Control and Traceability on
                Hedera}, with fellow ambassador Juan Manuel G\'omez.
            \item Judged on behalf of Hedera at ETHOnline 2026 (about \textbf{200 projects}) and
                NEA Tech 2026, Argentina (about \textbf{20 projects}).
            \item Onboard new developers and support builders on smart contracts and SDKs;
                represented Hedera at ETHDenver 2026, Hedera DevDay and Consensus Miami 2026.
        \end{itemize}

    \item \entry{Technical Lead}{X-Culture Professional Association (XPA)}{2025 -- 2026}
        \begin{itemize}
            \item Built the official website (WordPress, Elementor) with payment integration and
                ran the technical side of international virtual events.
        \end{itemize}

    \item \entry{Programming Instructor \& Mentor}{Kodland}{2021 -- 2024}
        \begin{itemize}
            \item Taught \textbf{Python} to \textbf{150+ students} (ages 10 to 20) in 30+ courses.
        \end{itemize}
\end{itemize}

\needspace{6\baselineskip}
\section*{Featured Projects}
\begin{itemize}
    \needspace{4\baselineskip}
    \item \textbf{Tacaño} (\gh{tacano}), open-source Claude Code plugin: cuts token use by about
        \textbf{39\%}, measured by replaying \textbf{500 real turns}. A hook asks a small model
        when to compact and what to hand to cheaper workers. Listed in a plugin marketplace,
        docs in English and Spanish.
    \needspace{4\baselineskip}
    \item \textbf{Willys por Colombia} (\href{https://willys-por-colombia.vercel.app}{live demo},
        \gh{willys-por-colombia}): driving game in three.js scored live by \textbf{Google Lyria
        RealTime}. Music mixes by distance to each region and drifts as you drive; each region
        replays its theme from a fixed seed. 18 automated test suites, touch controls.
    \needspace{4\baselineskip}
    \item \textbf{VerifiChain Visual} (\gh{verifichain-vision-oracle}), Hedera APEX 2026
        hackathon: visual inspection agent for metal and concrete (YOLOv8 with a Claude
        fallback). Each verdict is recorded on Hedera and certified with an NFT. Built around
        the needs of industrial companies in Santander.
    \needspace{4\baselineskip}
    \item \textbf{x402 Guard} (\gh{x402-guard}), Hedera x402 Bounty: payment proxy for AI
        agents. The agent holds no keys; the proxy pays only when a spending policy allows it,
        and every decision, approved or blocked, goes to a public audit trail (Hedera
        Consensus Service).
    \needspace{4\baselineskip}
    \item \textbf{Hive Protocol} (\gh{HIVE}), ETHDenver 2026, \textbf{2nd place} in the Hedera
        on-chain automation bounty: three AI agents propose, execute and validate knowledge
        with no human in the loop, using an MCP orchestrator and on-chain topics.
    \item \textbf{More:} hackathon projects ScoutBrief, Quorum Treasury, PULSO Tanda and Roast
        Pitch (\href{https://github.com/Nicolas6879}{github.com/Nicolas6879}).
\end{itemize}

\needspace{6\baselineskip}
\section*{Education}
\begin{itemize}
    \item \entry{Universidad Industrial de Santander (UIS)}{B.S. in Systems Engineering}{2021 -- Jul.\ 2026}
    \item \entry{SENA}{Technician in Industrial Electronic Equipment}{2018 -- 2019}
\end{itemize}

\needspace{6\baselineskip}
\section*{Technical Skills}
\begin{itemize}
    \item \textbf{Languages:} Python, TypeScript, JavaScript, PHP, Visual Basic, SQL, Solidity
        (also Rust, Java, C++, Go)
    \item \textbf{Web:} Vue 3, React, Next.js, three.js, Laravel, Node.js, REST
    \item \textbf{AI \& Agents:} MCP, Claude Code plugins and hooks, tool calling, multi-agent
        systems, LLM APIs (Claude, Gemini, OpenAI), YOLOv8
    \item \textbf{Blockchain:} Hedera (consensus service, tokens, smart contracts), Foundry, Hardhat
    \item \textbf{Data \& Tools:} SQL Server, MySQL, PostgreSQL, Power BI, Docker, Git, Vercel, n8n
    \item \textbf{Reverse engineering:} traffic analysis, Android and Flutter binaries, REA
\end{itemize}

\section*{Languages}
Spanish (native) \textperiodcentered{} English (intermediate, B1; improving toward B2)

\end{document}
